% Generated from completed homotaurine training and independent analytical reference.
\subsection{A second molecular application: homotaurine}
The same normalization and four-state secular model were applied to the
supplied illustrative homotaurine vector
$(J_A,J_M,J_X,d,e)=(-12.743923,-13.969277,-14.109177,3.093406,2.344390)$ Hz.
The temperature label of this historical fit is unresolved (289/298 K);
these calculations are not presented as a new experimental refit.
The normalized inputs are $(x_A,x_X,u,v)=((J_A-J_M)/s,(J_X-J_M)/s,d/s,e/s)$,
where $s=-J_M>0$ and $E_i=(-3J_A+J_M-3J_X)/4+s\lambda_i$.
Three unequal geminals and type-II mixing are retained.

As in the DSS experiment, a (4,16,4) ASN with $\{0,x,x^2\}$ activations
was trained on 100000 numerically diagonalized Hamiltonians, with disjoint
80000/10000/10000 training/validation/test splits and five initialization seeds.
Each geminal varied independently by $\pm0.3$ Hz and each coupling difference
by $\pm20\%$ about homotaurine. Training used Adam ($10^{-3}$), MAE,
batch size 4096, 900 epochs of softmax with $0.1L_{0.5}$ and 100 epochs of
sparsemax without the penalty. Input/output standardization used only training
statistics and was undone algebraically in the extracted expressions.
Validation eigenvalue RMSE alone selected the checkpoint and representative seed.
Eigenvalue order was preserved in the target labels; predictions were not sorted afterward.

The selected local ASN (seed 45) achieved frequency RMSE 0.002135162 Hz and maximum line error 0.02108312 Hz on 10000 held-out samples, with 0 ordering violations. Quadratic and quartic least-squares RMSEs on the same data were 0.002049504 and 5.946765e-05 Hz. At the supplied homotaurine reference, the analytical frequencies were 1.967923, 2.185440, 4.153363 Hz and the ASN frequencies were 1.968174, 2.185389, 4.153563 Hz. Three-decimal coefficient rounding gave a local frequency RMSE of 0.01304336 Hz.

With geminals fixed and $d=3.093406t$, $e=2.344390t$, the analytical
frequencies follow from the ordered roots of $z^3-C_1z^2+C_2z-C_3=0$, where
\[C_1=2.852388810881+0.077202373264t^2,\]
\[C_2=2.701400997954+0.073465038627t^2+0.001381140469t^4,\quad C_3=0.849012187073.\]
Writing $\omega_i=s\sqrt{z_i}$ gives $f_{12}=\omega_3-\omega_2$,
$f_{23}=\omega_2-\omega_1$, and $f_{13}=f_{12}+f_{23}$.
Substitution into the trained four-input ASN, without an additional section fit, gives
\[\widehat f_{12}=- 0.062329221 t^{2} + 1.904871944 t + 0.12563094\quad\mathrm{Hz}.\]
\[\widehat f_{23}=0.281058571 t^{2} + 1.251177692 t + 0.653152603\quad\mathrm{Hz}.\]
\[\widehat f_{13}=0.21872935 t^{2} + 3.156049636 t + 0.778783543\quad\mathrm{Hz}.\]
These are finite-domain learned approximations for $0.8\le t\le1.2$,
not exact perturbative Taylor coefficients. Spectral overlays use common
illustrative local-singlet amplitudes and 0.1 Hz Lorentzian FWHM, changing only
the line centers between exact and learned predictions.

Dropping type-II mixing shifts a principal line by 0.053710 Hz; replacing
geminals by their mean shifts a line by 0.278275 Hz. The second-order effective
three-state Hamiltonian has a maximum line error of 0.000287194 Hz at this point.
The fitted end-to-end difference $q=0.027306$ Hz changes a line by up to
0.018095 Hz and produces matched-parity separations up to 0.036090 Hz.
It should be included in precision comparisons with measured spectra; it is
omitted explicitly in both the analytical and learned four-input benchmark.
Full parity-resolved calculations and approximation errors accompany the comparison.

These results demonstrate explicit learned molecular approximations in a
specified domain. Polynomial and effective-Hamiltonian baselines are retained;
no intrinsic accuracy advantage over an appropriate fixed basis is claimed.
The broad-box run, rounded-coefficient errors, seed variation and decoupled-boundary
ordering checks are reported in the supplementary material.
